\documentclass[bibliography=totocnumbered]{scrartcl}
\input{settings/packages}
\input{settings/commands}
\author{Autor 1\\(Matrikelnummer)\\Mail \and Autor 2\\(Matrikelnummer)\\Mail}

\hypersetup{
	colorlinks=true,
	linkcolor=blue,
	filecolor=magenta,      
	urlcolor=cyan,
	citecolor=lime!50!black,
	filecolor=red
}
%\addbibresource{} %Bibliographiedateien laden
\addbibresource{Quellen.bib}

\geometry{a4paper, left=25mm, right=25mm, top=25mm, bottom=25mm}
\title{Titel}
% Alternative titles
% - Metric Analysis for Crystalls using Density of States Plots
\date{\today}
\lhead{\thedate}
\lhead{\thetitle}
% \pagestyle{fancy}

\usetikzlibrary{patterns}
\usetikzlibrary{3d}
\pagestyle{plain}
\allowdisplaybreaks
% function for multiple white spaces
\newcommand{\ws}[1]{% ws steht für white spaces
  \ifnum#1>0
    \foreach \i in {1,...,#1}{\;}%
  \fi
}

\definecolor{codegreen}{rgb}{0,0.6,0}
\definecolor{codegray}{rgb}{0.5,0.5,0.5}
\definecolor{codepurple}{rgb}{0.58,0,0.82}
\definecolor{backcolour}{rgb}{0.95,0.95,0.92}

\lstdefinestyle{mystyle}{
    backgroundcolor=\color{backcolour},   
    commentstyle=\color{codegreen},
    keywordstyle=\color{magenta},
    numberstyle=\tiny\color{codegray},
    stringstyle=\color{codepurple},
    basicstyle=\ttfamily\footnotesize,
    breakatwhitespace=false,         
    breaklines=true,                 
    captionpos=b,                    
    keepspaces=true,                 
    numbers=left,                    
    numbersep=5pt,                  
    showspaces=false,                
    showstringspaces=false,
    showtabs=false,                  
    tabsize=2
}

\lstset{style=mystyle}

\AtBeginBibliography{%
  \raggedright
  \sloppy
}

\setlength{\LTpre}{0pt}
\setlength{\LTpost}{0pt}
\title{\Huge Titel\\
    \huge \ac{GPR} Nr.
}

\rfoot{Versuchsnummer}

\begin{document}
    \renewcommand{\abstractname}{Abstract} % Nur zur verwenden, wenn man nicht in englisch schriebt
    \begin{titlepage}
        \centering
        \maketitle
        \vspace{1cm}
        \begin{tabular}{ll}
				\textbf{Betreuer:} & Name \\ 
				& Mail-Adresse\\
				\textbf{Raum:} &  \\ 
				\textbf{Messplatz:} &  \\
		\end{tabular}
        \vspace{1cm}
        \begin{abstract}
            \noindent Das Template kann gerne kopiert werden. Wichtig ist dabei, dass man den Settings-Ordner mitnimmt, damit das \TeX-Dokument funktioniert. Es wird davon abgeraden, die Inhalte vom GPR1/GPR2 zu kopieren (wegen Plagiard). Zudem sollten ein, zwei Ergebnisse aus der Analyse hier in den Abstract (Zusammenfassung mit eingebunden werden.
            \begin{table}[H]
                \centering
                \begin{tabular}{c|c}
                   Ergebnis 1  & Ergebnis 2 \\
                   \hline
                    x1 & x2
                \end{tabular}
                % \caption{Caption}
                % \label{tab:my_label}
            \end{table}
        \end{abstract}
    \end{titlepage}
    \tableofcontents
    \listoffigures
    \listoftables
    \input{settings/acronyms_de}
    \newpage
    \section{Einführung}
        \begin{itemize}
            \item Warum ist der Versuch hilfreich?
            \item Existiert \ac{bspw.} ein oder zwei Paper auf Google Scholar, die eine Antwort auf die vorherige Frage liefern
        \end{itemize}
    \section{Theorie und Versuchsaufbau}
        Welche Formeln wurden verwenden und warum? Wie sieht der Versuchsaufbau aus? 
        \begin{figure}
            \centering
            \includegraphics[width=\linewidth]{F}
            \caption{Caption}
            \label{fig:enter-label}
        \end{figure}
    \section{Analyse}
        Wie sehen die Ergebnisse der Analyse aus? Wie können diese interpretiert werden?
    \section{Diskussion}
        Diskussion der Ergebnisse und Vergleich mit anderen Daten (andere Protokolle, Paper, Bücher, \ac{etc.}), mit Angabe der Quellen.

    \newpage
    \appendix
    \section{Anhang}
    Versuchsaufzeichnungen
    \newpage
    \printbibliography[title={Literatur}]
\end{document}